\section{Mathematically Grounded Metrics for Data Quality}

To operationalize the Staged Gate Protocol, we must define specific metrics. These metrics extend beyond simple error checking; they quantify the \textit{fitness} of the data for the specific task of causal physics-informed learning.

\subsection{Symplectic Conservation Score (SCS)}

While Brutus is an arbitrary-precision integrator, it is not inherently symplectic (it does not strictly preserve phase space volume by construction, unlike Leapfrog). However, for the data to be valid training material for a physics-informed model, it must exhibit the properties of a conservative Hamiltonian system $H(\vect{q}, \vect{p})$.

\paragraph{Derivation:}
The Hamiltonian for the N-body problem is:
\begin{equation}
    H = \sum_{i} \frac{\vect{p}_i^2}{2m_i} - \sum_{i < j} \frac{G m_i m_j}{|\vect{q}_i - \vect{q}_j|}
\end{equation}
For a physically valid solution, $\frac{dH}{dt} = 0$. In numerical data, errors manifest as energy drift.

\paragraph{Metric Definition:}
We distinguish the raw trajectory-level Hamiltonian drift from its logarithmic score. The raw drift over the trajectory duration $T$ is:
\begin{equation}
    \epsilon_H = \max_{t \in [0, T]} \left| \frac{H(t) - H(0)}{H(0)} \right|
\end{equation}
The Symplectic Conservation Score ($SCS$) is the base-10 logarithm of this raw drift:
\begin{equation}
    SCS = \log_{10}(\epsilon_H)
\end{equation}
Throughout this document, $\epsilon_H$ denotes the raw Hamiltonian drift; $SCS$ denotes its base-10 logarithm. If $H(0)$ is zero or numerically near zero, $SCS$ is undefined; flag the trajectory for manual review or resimulation rather than forcing it through the logarithmic gate.

\paragraph{Thresholds for CTL-PINN:}
\begin{itemize}
    \item \textbf{$SCS \leq -10$ ($\epsilon_H \leq 10^{-10}$):} Excellent. Brutus has converged to high precision; ideal for training.
    \item \textbf{$-10 < SCS \leq -6$ ($10^{-10} < \epsilon_H \leq 10^{-6}$):} Acceptable. The drift is small enough that the trajectory data and Hamiltonian-conservation loss should remain mutually consistent during training.
    \item \textbf{$-6 < SCS \leq -5$ ($10^{-6} < \epsilon_H \leq 10^{-5}$):} Marginal. Keep only for review, diagnostics, or resampling decisions; do not use as clean training data without explicit justification.
    \item \textbf{$SCS > -5$ ($\epsilon_H > 10^{-5}$):} Reject. The trajectory exhibits significant Hamiltonian drift. A PINN trained on this will face conflicting gradients: the data loss tries to fit the drifting trajectory while the physics loss tries to enforce $\dot{H}=0$, leading to convergence failure.
\end{itemize}

\paragraph{Boundary Behavior:}
The original thresholds used strict excellent, acceptable, and reject bands around $-10$, $-6$, and $-4$, leaving $SCS = -10$, $SCS = -6$, and the interval $[-6, -4]$ underdefined. The corrected thresholds close these gaps with non-strict boundaries and introduce a Marginal band for $-6 < SCS \leq -5$.

\textbf{Insight:} In chaotic systems, energy error often grows linearly or quadratically with time in non-symplectic integrators. A high SCS suggests the integrator step size was too large for the chaotic amplification factor of the specific orbit.

\subsection{Stiffness-Aware Sampling Density (SASD)}

Stiffness is the primary cause of failure for neural ODE solvers and PINNs. It occurs when the timescales of the system's components vary by orders of magnitude—common in N-body problems during close encounters where the dynamical time $\tau \sim \sqrt{r^3/GM}$ shrinks rapidly. Brutus handles this adaptively, but if the output data is saved at fixed intervals (or if the adaptive steps are too coarse for the \textit{network's} capacity), the data will be ``aliased'' relative to the physics.

\paragraph{Metric Derivation:}
We define a local stiffness proxy based on the curvature of the trajectory, which correlates with the acceleration magnitude $\norm{\vect{a}}$. The sampling density is $\rho(t) = 1/\Delta t_i$.
The metric compares the time step $\Delta t_i$ to the local characteristic timescale $\tau_{char} = \frac{\norm{\vect{v}}}{\norm{\vect{a}}}$ (time to change velocity direction significantly).

\paragraph{Metric Definition:}
\begin{equation}
    SASD_i = \frac{\Delta t_i}{\tau_{char}(t_i)} = \Delta t_i \cdot \frac{\norm{\vect{a}(t_i)}}{\norm{\vect{v}(t_i)}}
\end{equation}
(To avoid division by zero at turning points, a more robust form uses the radius of curvature $\rho_{curv} = v^2/a_{\perp}$).

\paragraph{Interpretation:}
\begin{itemize}
    \item \textbf{$SASD \ll 1$:} The trajectory is over-sampled. The path between points is effectively linear. Safe for training.
    \item \textbf{$SASD \sim 1$:} The trajectory turns significantly between data points.
    \item \textbf{$SASD > 1$:} The trajectory is under-sampled. The particle performs a ``sharp turn'' or ``cusp'' between data points.
\end{itemize}

\textbf{Risk to CTL-PINN:} A high SASD value implies that the ``causal link'' between $t_i$ and $t_{i+1}$ is broken. The physics residual $\mathcal{L}_{physics}$ evaluated at a collocation point \textit{between} these data points (using interpolation) will be massive because the interpolant cannot capture the cusp. This stops the propagation of causal weights. Data with $SASD > 0.5$ should be flagged for resampling.

\subsection{Interpolation Physics Residual (IPR)}

This is a critical ``second-order'' metric. PINNs do not just learn the data points; they learn a continuous manifold. Therefore, the \textit{interpolant} implied by the data must satisfy the physics. If we use Cubic Hermite Splines (which preserve positions and velocities) to interpolate the Brutus data, does this spline satisfy Newton's laws?

\paragraph{Protocol:}
\begin{enumerate}
    \item Construct a piecewise Cubic Hermite Interpolating Polynomial (PCHIP) $\vect{q}_{spline}(t)$ using the data $\{(\vect{q}_i, \vect{v}_i)\}$.
    \item Sample $M$ synthetic ``collocation points'' $\tau_k$ uniformly distributed \textit{between} the data points.
    \item Calculate the physics residual at these synthetic points:
    \begin{equation}
        IPR = \frac{1}{M} \sum_{k=1}^{M} \left\| \frac{d^2 \vect{q}_{spline}}{dt^2}(\tau_k) - \vect{F}_{gravity}(\vect{q}_{spline}(\tau_k)) \right\|^2
    \end{equation}
\end{enumerate}

\paragraph{Relevance:}
If $IPR$ is high, it means that even though the data points are correct, the ``shape'' of the trajectory between them (as defined by standard interpolation) is unphysical. This indicates that the data density is insufficient to constrain the continuous solution space of the PINN. The network will struggle to find a function that fits both the data points and the physics equation between them, leading to ``wobbling'' or high-frequency artifacts. This metric directly predicts the ``Physics Loss'' component of the PINN training.

\subsection{Phase Space Coverage Entropy (PSC)}

Chaotic systems are ergodic (or have ergodic components). For a model to generalize, the training data must cover the relevant phase space structures—both the stable regions (quasi-periodic orbits) and the unstable regions (close encounters, scattering). A dataset biased towards one will fail on the other.

\paragraph{Metric Definition:}
We discretize the phase space of interest (e.g., the $r_{12}, r_{23}$ distance plane or the Poincaré section) into a grid of $K$ bins. Let $p_k$ be the probability of a data point falling into bin $k$. The Shannon Entropy is:
\begin{equation}
    H_{data} = - \sum_{k=1}^{K} p_k \log p_k
\end{equation}

\paragraph{Application:}
\begin{itemize}
    \item \textbf{Low Entropy:} The data is clustered in a small region (e.g., all particles far apart). The model will fail to learn close-encounter dynamics.
    \item \textbf{High Entropy:} The data explores diverse configurations.
    \item \textbf{Stiffness Coverage Ratio:} Specifically check the bin corresponding to $r_{min} \to 0$. If this bin is empty, the dataset lacks the ``hard'' physics examples necessary to train the network to handle singularities. The report recommends ``Importance Sampling'' to enrich the dataset with these high-stiffness events.
\end{itemize}
